\phantomsection\addcontentsline{toc}{mysection}{摘\ \ \ \ \ \ \ \ 要}
\section*{\zihao{-3}\heiti 光线追踪渲染器实现与基于BVH和重要性采样的优化}
\section*{\zihao{3}\heiti 摘\ \ \ \ \ \ \ \ 要}
\zihao{-4}光线追踪渲染作为数学与计算机图形学交叉学科中的重要技术，在实现逼真的图像渲染方面发挥着关键作用。此研究的意义在于其为视觉效果和图形应用的开发提供了强大的工具，能够在虚拟环境中创造出更加真实的视觉体验。本文通过模拟光线在显示世界中的传播与反射方式，基于C++语言构建了一个不依赖专用硬件的光线追踪渲染器，能够实现基本几何体、基本材质和简单纹理的渲染。结合计算机图形学、大数据和数值分析方法（如BVH、概率密度估计、重要性采样）对可能出现的噪点和采样数问题进行优化。最后通过统计和数字图像处理方法，对渲染结果进行评估，指出本文解决方案的优点与不足。深入研究光线追踪算法及其优化的数学方法，可以提升渲染效率和图像质量，为数学、大数据和计算机图形学相关领域与交叉学科的发展做出贡献。其研究和应用对于大数据领域正如火如荼发展的数字孪生技术也有重要意义，更加精确的克隆体不仅可以带来更好的用户视觉体验，更可以精准复制模拟物理世界的状态，有望成为大数据数字孪生视觉领域的技术底座。
\\
%空一行
\\
\zihao{-4}{\heiti 关键词：}计算机图形学；光线追踪；层次包围盒；蒙特卡洛积分；重要性采样

\clearpage

\phantomsection\addcontentsline{toc}{mysection}{ABSTRACT}
\section*{\zihao{-3}\bfseries Ray Tracing Renderer Implementation with Optimisation Based on BVH and Importance Sampling}

\section*{\zihao{3}\bfseries ABSTRACT}
\zihao{-4}Ray tracing rendering, as an important technique in the interdisciplinary field of mathematics and computer graphics, plays a key role in achieving realistic image rendering. The significance of this research lies in the fact that it provides a powerful tool for the development of visual effects and graphical applications, which can create a more realistic visual experience in virtual environments. In this paper, by simulating the way light propagates and reflects in the display world, we build a ray-tracing renderer based on C++ that does not rely on dedicated hardware and is capable of rendering basic geometries, basic materials and simple textures. Combining computer graphics, big data and numerical analysis methods (e.g. BVH, probability density estimation, importance sampling) to optimise possible noise and sample count problems. Finally, the rendering results are evaluated by means of statistics and digital image processing methods, pointing out the strengths and weaknesses of the solutions in this paper. An in-depth study of ray tracing algorithms and the mathematical methods for their optimisation can improve the rendering efficiency and image quality, contributing to the development of related fields and interdisciplinary disciplines of mathematics, big data and computer graphics. Its research and application is also of great significance to the digital twin technology that is developing in full swing in the field of big data, and more accurate clones can not only bring better user visual experience, but also accurately replicate the simulation of the state of the physical world, which is expected to become a technological pedestal for the field of big data digital twin vision.
\\
%空一行
\\
\zihao{-4}{\bfseries Key words:\ }Computer graphics; Ray tracing; Bounding volume hierarchies; Monte Carlo integration; Importance sampling

\clearpage




